\subsubsection{椭圆方程组的补条件}\label{subsec:elliptic-system-complementing-condition}

对于 ODE 边值问题, 众所周知, 若每个端点的边界条件数等于微分方程阶数的一半, 且每个端点的边界条件线性无关, 就得到一个 ``合理\textquotesingle\textquotesingle\ 的边值问题. 对椭圆 PDE, 情况更加复杂. 考虑问题
\begin{equation}\label{eq:complementing-motivation-biharmonic-equation}
\Delta\Delta u=0
\end{equation}
及边界条件
\begin{equation}\label{eq:complementing-motivation-biharmonic-boundary}
\Delta u=\frac{\partial}{\partial n}\Delta u=0.
\end{equation}
虽然这两个边界条件相互独立, 但每个调和函数同时满足微分方程与边界条件. 即使在多项式类中, 也存在无穷多个线性无关的调和函数. 这说明需要对边界条件提出能产生 ``好\textquotesingle\textquotesingle\ 问题的假设. 这些假设不仅表达边界条件之间的独立性, 还涉及边界条件与微分方程之间的关系.

补条件给出了对这种 ``好\textquotesingle\textquotesingle\ 边界条件的刻画. 将对一般椭圆方程组陈述它. 与 Chapter~\ref{chap:characteristics} 一样, 考虑一个 $k\times k$ 方程组
\begin{equation}\label{eq:weighted-elliptic-general-system}
L_{ij}(\bm{x},D)u_j=f_i(\bm{x}),\quad\bm{x}\in\Omega,\ i=1,\ldots,k.
\end{equation}
如前所述, 给第 $i$ 个方程赋予 ``权重\textquotesingle\textquotesingle\ $s_i$, 给第 $j$ 个自变量赋予权重 $t_j$, 使得 $L_{ij}$ 的阶数至多为 $s_i+t_j$, 并用 $L_{ij}^p$ 表示恰为 $s_i+t_j$ 阶的项. 椭圆性条件为
\begin{equation}\label{eq:weighted-elliptic-system-principal-determinant}
\det\bm{L}^p(\bm{x},\bm{\xi})\ne0\quad\forall\bm{\xi}\in\mathbb{R}^n\backslash\{0\}.
\end{equation}
如 Chapter~\ref{chap:characteristics} 中所述, 该条件可解释如下: 在 $\Omega$ 中任取固定点 $\bm{x}_0$, 取该点处的系数值, 并在整个 $\mathbb{R}^n$ 上考虑方程组 $\bm{L}^p(\bm{x}_0,D)\bm{u}=0$. 椭圆性意味着这个常系数方程组没有非常数的周期解. 补条件将是在边界点处的相应条件.

设方程组的阶数, 即多项式 $\det\bm{L}^p(\bm{x},\bm{\xi})$ 的阶数, 为 $2m$. 施加 $m$ 个边界条件
\begin{equation}\label{eq:weighted-elliptic-general-boundary-conditions}
B_{lj}(\bm{x},D)u_j=g_l(\bm{x}),\quad\bm{x}\in\partial\Omega,\ l=1,\ldots,m.
\end{equation}
现在给每个边界条件 $l=1,\ldots,m$ 定义权重 $r_l$, 使 $B_{lj}$ 的阶数不超过 $r_l+t_j$. 同样, 恰为 $r_l+t_j$ 阶的项视为主部. (若 $r_l+t_j$ 为负, 当然理解为 $B_{lj}=0$.) 现在设 $\bm{x}_0$ 为 $\partial\Omega$ 上的一点, $\bm{n}$ 为 $\Omega$ 的外法向量. 在半空间 $(\bm{x}-\bm{x}_0)\cdot\bm{n}<0$ 上考虑常系数问题
\begin{equation}\label{eq:complementing-frozen-principal-system}
L_{ij}^p(\bm{x}_0,D)u_j=0,\quad i=1,\ldots,k,
\end{equation}
并在边界 $(\bm{x}-\bm{x}_0)\cdot\bm{n}=0$ 上施加边界条件
\begin{equation}\label{eq:complementing-frozen-principal-boundary}
B_{lj}^p(\bm{x}_0,D)u_j=0,\quad l=1,\ldots,m,
\end{equation}


\begin{Definition}
\label{def:elliptic-system-complementing-condition}
若 \eqref{eq:complementing-frozen-principal-system}, \eqref{eq:complementing-frozen-principal-boundary} 没有以下形式的非平凡解, 则称\textbf{补条件}在 $\bm{x}_0$ 处成立:
\begin{equation}\label{eq:complementing-decaying-tangential-mode}
\bm{u}(\bm{x})=\exp(i\bm{\xi}\cdot(\bm{x}-\bm{x}_0))\bm{v}(\eta),
\end{equation}
其中 $\bm{\xi}$ 为垂直于 $\bm{n}$ 的非零实向量, $\eta=(\bm{x}-\bm{x}_0)\cdot\bm{n}$, 且当 $\eta\to-\infty$ 时, $\bm{v}(\eta)$ 指数趋于 $0$.
\end{Definition}
